\providecommand{\Res}[2]{\qty[round-mode=figures,round-precision=3]{\fpeval{#1}}{#2}}
\providecommand{\Val}[1]{\num[round-mode=figures,round-precision=4]{\fpeval{#1}}}
\providecommand{\ResSci}[2]{\qty[round-mode=figures,round-precision=3,exponent-mode=scientific]{\fpeval{#1}}{#2}}
\subsubsection*{Question a) Décalage}
\begin{enumerate}[label=\arabic*)]
    \item \textbf{Variables connues et recherchées :} $\theta = \qty{\v{th}}{\degree}$, $h/(m_e c) = \qty{2.425e-12}{m}$ ; variable recherchée : $\Delta\lambda$
    \item \textbf{Équations pertinentes :}
    \[ \Delta\lambda = \frac{h}{m_e c}(1 - \cos\theta) \]
    \item \textbf{Isolation de la variable recherchée :} l'équation est déjà isolée.
    \item \textbf{Résolution numérique :}
    \[ \Delta\lambda = \num{2.425e-12} \cdot (1 - \cos \v{th}^\circ) = \ResSci{6.626e-34/(9.109e-31*3e8)*(1 - cosd(\v{th}))}{m} \]
    \item \textbf{Réponse :} $\Delta\lambda = \mathbf{\Res{1e12*6.626e-34/(9.109e-31*3e8)*(1 - cosd(\v{th}))}{pm}}$, le même pour les deux types de photons : il ne dépend que de l'angle.
\end{enumerate}

\subsubsection*{Question b) Lumière visible}
\[ \frac{\Delta\lambda}{\lambda} = \frac{\ResSci{6.626e-34/(9.109e-31*3e8)*(1 - cosd(\v{th}))}{}}{\num{\vsi{lv}}} = \ResSci{100*6.626e-34/(9.109e-31*3e8)*(1 - cosd(\v{th}))/\vsi{lv}}{\percent} \]

\subsubsection*{Question c) Rayons X}
\[ \frac{\Delta\lambda}{\lambda} = \frac{\ResSci{6.626e-34/(9.109e-31*3e8)*(1 - cosd(\v{th}))}{}}{\num{\vsi{lx}}} = \Res{100*6.626e-34/(9.109e-31*3e8)*(1 - cosd(\v{th}))/\vsi{lx}}{\percent} \]
Pour la lumière visible, la variation est d'environ un millionième de la longueur d'onde, impossible à mesurer. Pour les rayons X, la longueur d'onde est comparable à $h/(m_e c)$ et la variation de quelques pour cent est facilement mesurable.
